\section{Generative Models via Transportation}
\label{sec-generative-models-transportation}
\subsection{Generative Models via Flow Matching}
\label{sec-generative-flow-matching}
\paragraph{Deterministic interpolants.}
The word ``stochastic'' can hide two different levels of randomness. We first use the simpler one: after drawing a latent variable $U\sim\pi$, the path $t\mapsto I_t(U)$ is deterministic and differentiable.

\begin{defn}[Deterministic interpolant]\label{def-deterministic-interpolant}
Let \(\pi\in\Couplings(\alpha_0,\alpha_1)\). A deterministic interpolant is a measurable family \(I_t:\RR^d\times\RR^d\to\RR^d\), differentiable in $t$, satisfying
\(I_0(x,y)=x, \qquad I_1(x,y)=y.\)
It induces \(\alpha_t=(I_t)_\sharp\pi\) and has pathwise velocity \(\partial_t I_t(x,y)\).
\end{defn}
\begin{equation}\label{eq:interp-coupling}
\forall t \in [0,1], \quad \alpha_t:= (I_t)_\sharp \pi.
\end{equation}
If $\pi = \alpha \otimes \beta$ and $\alpha = \frac{1}{n} \sum_i \delta_{x_i}$, $\beta = \frac{1}{m} \sum_j \delta_{y_j}$, then $\alpha_t$ consists of $n \times m$ Dirac masses

\(\alpha_t = \frac{1}{nm} \sum_{i,j} \delta_{I_t(x_i,y_j)}.\)
If $\pi = (\Id, T)_\sharp \alpha$ is a Brenier-type coupling, then $\alpha_t = ((1-t)\Id + tT)_\sharp \alpha$ is the so-called McCann OT interpolation.

\paragraph{Flow matching formula.}
\begin{prop}[Flow matching vector field]\label{prop-flow-matching-vector-field}
For each fixed $t$, assume $\partial_t I_t\in L^2(\pi;\RR^d)$. Consider the flow-matching problem over measurable fields $v_t:\RR^d\to\RR^d$
\begin{equation}
\min_{v_t} \int_{\RR^{d'}} \norm{v_t(I_t(u)) - [\partial_t I_t](u)}^2 \, \d\pi(u). \label{eq:flow-matching}
\end{equation}
Its minimizer is characterized $\alpha_t$-almost everywhere by the conditional expectation
\begin{equation}\label{eq:flow-match-conditional}
v_t(z) = \EE_{u \sim \pi} \big( [\partial_t I_t](u) \, \big| \, z = I_t(u) \big).
\end{equation}
Then the pair $(\alpha_t,v_t)$ satisfies the continuity equation~\eqref{eq:eulerian-advection}.
\end{prop}
\begin{proof}
We first recall the two equivalent ways of writing the interpolated measure. Formally, one may write
\(\alpha_t(z)=\int_{\RR^{d'}}\delta(z-I_t(u))\,\d\pi(u),\)
while the rigorous meaning is that, for every smooth test function $\varphi$,
\begin{equation}\label{eq:flow-matching-pushforward-test}
\int_{\RR^d}\varphi(z)\,\d\alpha_t(z)
=
\int_{\RR^{d'}}\varphi(I_t(u))\,\d\pi(u).
\end{equation}
The minimizer in~\eqref{eq:flow-matching} is the orthogonal projection in $L^2(\pi;\RR^d)$ of the latent velocity $\partial_t I_t(u)$ onto the closed subspace of functions that depend on $u$ only through $I_t(u)$. This projection is the conditional expectation~\eqref{eq:flow-match-conditional}. Formally, this can be read as
\(v_t(z)=\frac{1}{\alpha_t(z)} \int_{\RR^{d'}}\delta(z-I_t(u))[\partial_t I_t](u)\,\d\pi(u),\)
and rigorously it means that, for every smooth test vector field $m$,
\begin{equation}\label{eq:v_t}
\int \dotp{m(z)}{v_t(z)} \, \d\alpha_t(z)
=
\int \dotp{m(I_t(u))}{[\partial_t I_t](u)} \, \d\pi(u).
\end{equation}

The weak form of
$\partial_t\alpha_t+\diverg(\alpha_t v_t)=0$ is that, for every smooth scalar test function $\varphi$,
\begin{equation}\label{eq:flow-matching-weak-target}
\frac{\d}{\d t}\int\varphi(z)\,\d\alpha_t(z)
-
\int\dotp{v_t(z)}{\nabla\varphi(z)}\,\d\alpha_t(z)
=0.
\end{equation}
Using~\eqref{eq:flow-matching-pushforward-test} and differentiating under the integral sign gives
\begin{equation}\label{eq:flow-matching-test-derivative}
\frac{\d}{\d t}\int \varphi(z)\d\alpha_t(z)
=
\int \dotp{\nabla\varphi(I_t(u))}{[\partial_t I_t](u)}\d\pi(u).
\end{equation}
On the other hand, applying~\eqref{eq:v_t} with $m=\nabla\varphi$ gives
\begin{equation}\label{eq:flow-matching-velocity-test}
\int\dotp{v_t(z)}{\nabla\varphi(z)}\,\d\alpha_t(z)
=
\int \dotp{\nabla\varphi(I_t(u))}{[\partial_t I_t](u)}\d\pi(u).
\end{equation}
Comparing~\eqref{eq:flow-matching-test-derivative} and~\eqref{eq:flow-matching-velocity-test} yields~\eqref{eq:flow-matching-weak-target}, which is the desired continuity equation.
\end{proof}
The conditional expectation in~\eqref{eq:flow-match-conditional} has a simple measure-theoretic meaning. Let $\alpha_t=(I_t)_\sharp\pi$ and define the vector-valued flux measure $\omega_t$ on $\RR^d$ by

\(\int_{\RR^d}\dotp{\psi(z)}{\d\omega_t(z)} \eqdef \int_{\RR^{d'}}\dotp{\psi(I_t(u))}{[\partial_t I_t](u)}\d\pi(u)\)
\(\d\omega_t(z)=v_t(z)\d\alpha_t(z), \qquad v_t=\frac{\d\omega_t}{\d\alpha_t}.\)
In the language of Lebesgue decomposition, $\omega_t$ has only an absolutely continuous part with respect to $\alpha_t$ and no singular part; the conditional expectation is precisely its density. This agrees with the flux notation used in the dynamic formulation of Section~\ref{sec-benamou-brenier-dynamic}.

\(v_t(z)=\int_{\{I_t(u)=z\}}[\partial_t I_t](u)\d\pi_{t,z}(u).\)
Thus the solution of \eqref{eq:flow-matching} is the conditional expectation of the velocities $\partial_t I_t$: intuitively, $v_t(z)$ is the average velocity of all trajectories passing through $z$.

For the exact field $v_t$, integrating the ODE $\dot{x}=v_t(x)$ defines a transport map $T_t$. If $v_t$ is regular enough, or more generally if the continuity equation has a unique solution for this velocity, then $(T_t)_\sharp\alpha_0=\alpha_t$.

\paragraph{Static-noise stochastic interpolants.}
\begin{defn}[Stochastic interpolant]\label{def-stochastic-interpolant}
Let \((X_0,X_1)\sim\pi\in\Couplings(\alpha_0,\alpha_1)\), and let \(Z\) be an auxiliary random variable independent of \((X_0,X_1)\). A stochastic interpolant is \(X_t = I_t(X_0,X_1,Z),\)
where $t\mapsto I_t(x_0,x_1,z)$ is differentiable and
\(I_0(x_0,x_1,z)=x_0, \qquad I_1(x_0,x_1,z)=x_1.\)
Its marginal curve is \(\alpha_t=\operatorname{Law}(X_t)\).
\end{defn}
\(X_t=a(t)X_0+b(t)X_1+\gamma(t)Z, \qquad \gamma(0)=\gamma(1)=0.\)
The noise $Z$ is static: conditionally on $(X_0,X_1,Z)$, the path $t\mapsto X_t$ is differentiable. Thus this construction is exactly the previous push-forward framework with latent variable $u=(X_0,X_1,Z)$, law $\pi=\operatorname{Law}(X_0,X_1,Z)$, and interpolant given by the displayed $I_t$.

\(v_t(x)=\EE\bigl[\partial_t I_t(X_0,X_1,Z)\mid X_t=x\bigr],\)
\paragraph{Connection with diffusion models.}\label{par-diffusion-model-connection}
In the special case where $I_t(x,y)=(1-t)x+ty$ is a linear interpolation and $\pi = \alpha \otimes \beta$, the curve $\alpha_t$ is a convolution of rescaled versions of $\alpha_0$ and $\alpha_1$. The flow-matching problem~\eqref{eq:flow-matching} becomes

\(\min_{(v_t)_t} \int_{\RR^{d} \times \RR^d} \norm{v_t( (1-t)x+t y ) - (y-x) }^2 \, \d\alpha_0(x) \d\alpha_1(y).\)
\begin{proposition}[Tweedie identity]\label{prop:Tweedie}
Let $W$ be a random vector in $\RR^{d}$ with law $\beta\in\Pp_1(\RR^d)$.
For $\sigma>0$, observe
\(Z \;=\; W + \sigma\,\varepsilon, \quad\text{where } \varepsilon \sim \Gaussian(0,\Id) \text{ is independent of } W.\)
Denote by $\rho_\sigma$ the smooth positive density of
\(\beta_\sigma \;=\; \beta * \Gaussian\bigl(0,\sigma^{2}\Id\bigr),\)
which is the law of $Z$.
Then the conditional mean admits the following everywhere-defined version: \(\EE\bigl[\,W \mid Z=z\bigr] \;=\; z \;+\;\sigma^{2}\,\nabla \log \rho_\sigma(z) \qquad\text{for all } z \in \RR^{d}.\)
\end{proposition}
\begin{proof}
Let $\varphi_\sigma$ be the $\Gaussian(0,\sigma^{2}\Id)$ density. Bayes' rule gives \(\EE[W\mid Z=z] = \frac{1}{\rho_\sigma(z)} \int_{\RR^{d}} w\, \varphi_\sigma(z-w)\,\d\beta(w).\)
Differentiating the Gaussian convolution under the integral sign and using
$
\nabla_z\varphi_\sigma(z-w)
= -\sigma^{-2}(z-w)\,\varphi_\sigma(z-w)
$
yields
\(\nabla\rho_\sigma(z) = \int \nabla_z\varphi_\sigma(z-w)\,\d\beta(w) = -\sigma^{-2}\Bigl(z-\EE[W\mid Z=z]\Bigr)\,\rho_\sigma(z).\)
Rearranging finishes the proof.
\end{proof}
\begin{proposition}[Gaussian-endpoint flow-matching field]\label{prop:flow}
Let $X\sim\alpha$ and $Y\sim\Gaussian(0,\Id)$ be independent.
For $t\in(0,1)$ set
\(Z_t \;=\; (1-t)\,X + t\,Y, \qquad \alpha_t =\operatorname{Law}(Z_t)=\rho_t\,\d x.\)
The regression minimizer $v^\star:\RR^d\times(0,1)\to\RR^d$ of
\(\min_{v}\;\int_{0}^{1}\! \iint_{\RR^{d}\times\RR^{d}} \bigl|y-x-v\bigl((1-t)x+t y,t\bigr)\bigr|^{2}\, \d\alpha(x)\,\d\Gaussian(y)\,\d t\)
is \(v^\star(x,t) = -\frac{1}{1-t}\,x \;-\; \frac{t}{1-t}\,\nabla\log\rho_t(x) \qquad (x\in\RR^{d},\;t\in(0,1)).\)
In particular, for each $t\in(0,1)$ this field is a gradient field,
\[
v^\star(\cdot,t)=-\nabla
\left(
\frac{\norm{\cdot}^2}{2(1-t)}
+\frac{t}{1-t}\log\rho_t
\right).
\]
\end{proposition}
\begin{proof}
Fix $t\in(0,1)$ and write $W=(1-t)X$, $\sigma=t$, so that
$Z_t = W + \sigma\,Y$ matches the setting of Proposition~\ref{prop:Tweedie}.
Conditional expectations satisfy
$
v^\star(z,t)
= \EE[Y-X\mid Z_t=z]
= \frac{1}{t}\,\EE[Z_t-W\mid Z_t=z]
-\,\frac{1}{1-t}\,\EE[W\mid Z_t=z].
$
Applying Proposition~\ref{prop:Tweedie} to $\EE[W\mid Z_t=z]$ and
noting $\EE[Y\mid Z_t=z]
= -\,t\,\nabla\log\rho_t(z)$
gives the claimed formula.
\end{proof}
\paragraph{When is the induced map optimal?}
Integrating the learned velocity gives a deterministic map from $\alpha_0$ to $\alpha_1$, but this map is not automatically the Brenier optimal map. It is optimal only in special cases where the accumulated flow remains the gradient of a convex potential.

\begin{prop}[Gaussian flow matching and optimality]\label{prop-gaussian-flow-matching-optimality}
Let $\Sigma_0,\Sigma_1\succ0$ and let $X_0\sim\Gaussian(0,\Sigma_0)$ and $X_1\sim\Gaussian(0,\Sigma_1)$ be independent. Consider the linear flow-matching interpolation
\(Z_t=(1-t)X_0+tX_1, \qquad \alpha_t=\operatorname{Law}(Z_t)=\Gaussian(0,\Sigma_t),\)
\eqllead{where}{eq-gaussian-product-fm-covariance}{
\Sigma_t=(1-t)^2\Sigma_0+t^2\Sigma_1.
}
Then the exact flow-matching velocity is affine, $v_t(z)=A_tz$, with
\begin{equation}\label{eq-gaussian-product-fm-velocity}
A_t=\bigl(t\Sigma_1-(1-t)\Sigma_0\bigr)\Sigma_t^{-1}.
\end{equation}
The induced flow map $T_t^{\rm FM}$ from $\alpha_0$ to $\alpha_t$ is
\begin{equation}\label{eq-gaussian-product-fm-map}
T_t^{\rm FM}
=
\Sigma_0^{1/2}
\Bigl((1-t)^2\Id+t^2\Sigma_0^{-1/2}\Sigma_1\Sigma_0^{-1/2}\Bigr)^{1/2}
\Sigma_0^{-1/2}.
\end{equation}
\eqllead{In particular,}{eq-gaussian-product-fm-terminal-map}{
T_1^{\rm FM}
=
\Sigma_0^{1/2}
\bigl(\Sigma_0^{-1/2}\Sigma_1\Sigma_0^{-1/2}\bigr)^{1/2}
\Sigma_0^{-1/2}.
}
This terminal map coincides with the quadratic optimal transport map
\begin{equation}\label{eq-gaussian-brenier-map-comparison}
T^{\rm OT}
=
\Sigma_0^{-1/2}
\bigl(\Sigma_0^{1/2}\Sigma_1\Sigma_0^{1/2}\bigr)^{1/2}
\Sigma_0^{-1/2}
\end{equation}
if and only if $\Sigma_0\Sigma_1=\Sigma_1\Sigma_0$.
\end{prop}
\begin{proof}
The conditional-expectation formula gives
\(v_t(z)=\EE[X_1-X_0\mid Z_t=z]\).
Since all variables are jointly Gaussian, this conditional expectation is linear and
\(v_t(z) = \operatorname{Cov}(X_1-X_0,Z_t)\operatorname{Cov}(Z_t)^{-1}z = \bigl(t\Sigma_1-(1-t)\Sigma_0\bigr)\Sigma_t^{-1}z,\)
which proves~\eqref{eq-gaussian-product-fm-velocity}. To solve the characteristic equation, whiten the source by setting
\(C=\Sigma_0^{-1/2}\Sigma_1\Sigma_0^{-1/2}, \qquad \widetilde Z_t=\Sigma_0^{-1/2}Z_t.\)
In these coordinates the source covariance is $\Id$ and
\(\widetilde\Sigma_t=(1-t)^2\Id+t^2C\).
Because $\Id$ and $C$ commute, the affine flow map in whitened coordinates is simply $\widetilde T_t=\widetilde\Sigma_t^{1/2}$. Indeed,
\(\frac{\d}{\d t}\widetilde\Sigma_t^{1/2} = \bigl(tC-(1-t)\Id\bigr)\widetilde\Sigma_t^{-1/2},\)
which is exactly the equation $\dot{\widetilde T}_t=\widetilde A_t\widetilde T_t$ with $\widetilde T_0=\Id$. Returning to the original coordinates gives~\eqref{eq-gaussian-product-fm-map}, and $t=1$ gives~\eqref{eq-gaussian-product-fm-terminal-map}.

Both $T_1^{\rm FM}$ and $T^{\rm OT}$ push $\Gaussian(0,\Sigma_0)$ to $\Gaussian(0,\Sigma_1)$. The Brenier map between nondegenerate Gaussians is the unique symmetric positive definite linear map with this property. Hence $T_1^{\rm FM}=T^{\rm OT}$ if and only if $T_1^{\rm FM}$ is symmetric positive definite. The map $T_1^{\rm FM}$ is similar to $C^{1/2}$, so if it is symmetric then it is automatically positive definite. Since $C^{1/2}$ is symmetric positive definite,
\((T_1^{\rm FM})^\top = \Sigma_0^{-1/2}C^{1/2}\Sigma_0^{1/2}.\)
Thus symmetry of $T_1^{\rm FM}$ is equivalent to
$\Sigma_0 C^{1/2}=C^{1/2}\Sigma_0$, hence to $\Sigma_0 C=C\Sigma_0$ by functional calculus. Multiplying this identity on the left and right by $\Sigma_0^{1/2}$ gives $\Sigma_0\Sigma_1=\Sigma_1\Sigma_0$. Conversely, if $\Sigma_0$ and $\Sigma_1$ commute, they are orthogonally co-diagonalizable, and both~\eqref{eq-gaussian-product-fm-terminal-map} and~\eqref{eq-gaussian-brenier-map-comparison} reduce in that basis to the diagonal map with entries $\sqrt{\lambda_{1,k}/\lambda_{0,k}}$. This proves the equivalence.
\end{proof}
\paragraph{Variations on the interpolant.}
\(v_t(z)=\lambda'(t)\,v^{\rm lin}_{\lambda(t)}(z), \qquad v^{\rm lin}_{r}(z)=\EE[Y-X\mid (1-r)X+rY=z].\)
\(Z_t=a_tX+b_tY,\qquad Y\sim\Gaussian(0,\sigma^2\Id),\)
then both the mixture centers and the component variances are changed. Writing $p_t$ for the density of $Z_t$ and $s_t=\nabla\log p_t$, Tweedie's formula gives, away from times where $a_t=0$,

\[
v_t(z)=a'_t\,\EE[X\mid Z_t=z]+b'_t\,\EE[Y\mid Z_t=z]
=\frac{a'_t}{a_t}z+
\left(\frac{a'_tb_t^2}{a_t}-b'_tb_t\right)\sigma^2s_t(z).
\]
For the linear bridge, $a_t=1-t$ and $b_t=t$, this recovers the formula above. For the variance-preserving Ornstein--Uhlenbeck noising used in diffusion models,

\(a_\tau=e^{-\tau},\qquad b_\tau=\sqrt{1-e^{-2\tau}},\)
one obtains the forward probability-flow velocity $v_\tau(z)=-z-\sigma^2\nabla\log p_\tau(z)$. Sampling integrates this ODE backward in $\tau$.

\(a_t=(1-t)(1-2t), \qquad b_t=t,\)
\subsection{One-Step Generative Models}
\paragraph{One-step models via parameter-domain discrepancy flows.}
The most direct one-step strategy is to train the generator parameters themselves by descending a discrepancy between generated and data distributions. Let $\zeta$ be a simple latent distribution, let $f_\theta:\Zz\to\Xx$ be a neural generator, and let $\beta$ be the data distribution.

\((f_\theta)_\sharp\zeta=\beta,\)
or, in practice, such that the generated law is close to $\beta$. Given a discrepancy $\mathcal D$ on probability measures, define

\begin{equation}\label{eq-one-step-parameter-discrepancy}
\mathcal E_\beta(\gamma)\eqdef\mathcal D(\gamma,\beta),
\qquad
H_\beta(\theta)
\eqdef
\mathcal E_\beta\big((f_\theta)_\sharp\zeta\big)
=
\mathcal D\big((f_\theta)_\sharp\zeta,\beta\big).
\end{equation}
\begin{equation}\label{eq-one-step-parameter-flow}
\dot\theta_t=-\nabla_\theta H_\beta(\theta_t).
\end{equation}
\(\alpha_t\eqdef(f_{\theta_t})_\sharp\zeta, \qquad X_t=f_{\theta_t}(Z), \qquad Z\sim\zeta.\)
If $t\mapsto\theta_t$ is smooth and $f_\theta$ is differentiable with respect to $\theta$, then

\(\dot X_t = \partial_\theta f_{\theta_t}(Z)\dot\theta_t.\)
For $\alpha_t$-a.e. $x$, let $\eta_{t,x}$ be the disintegration of the latent law $\zeta$ with respect to the map $f_{\theta_t}$. Thus $\eta_{t,x}$ is supported on the fiber $\enscond{z}{f_{\theta_t}(z)=x}$, and for every bounded measurable $\psi$,

\[
\int \psi(z)\,\d\zeta(z)
=
\int_\Xx\left(\int \psi(z)\,\d\eta_{t,x}(z)\right)\d\alpha_t(x).
\]
\begin{equation}\label{eq-one-step-induced-velocity}
v_t(x)
=
\int \partial_\theta f_{\theta_t}(z)\dot\theta_t\,\d\eta_{t,x}(z)
=
-\int \partial_\theta f_{\theta_t}(z)\nabla_\theta H_\beta(\theta_t)\,\d\eta_{t,x}(z).
\end{equation}
\(\frac{\d}{\d t}\int \varphi(x)\,\d\alpha_t(x) = \int \dotp{\nabla\varphi(x)}{v_t(x)}\,\d\alpha_t(x),\)
\begin{equation}\label{eq-one-step-induced-continuity}
\partial_t\alpha_t+\diverg(\alpha_t v_t)=0.
\end{equation}
The velocity~\eqref{eq-one-step-induced-velocity} depends on the parameter value $\theta_t$ and on the latent disintegration, not only on the measure $\alpha_t$. If the parametrization is non-identifiable, two different parameters can generate the same measure while inducing different velocities.

\(-\nabla\delta_\gamma\mathcal E_\beta(\gamma)\big|_{\gamma=\alpha_t}.\)
\(\hat\beta_n=\frac{1}{n}\sum_{i=1}^n\de_{Y_i}, \qquad Y_i\overset{\mathrm{i.i.d.}}{\sim}\beta,\)
\[
\EE\left[\nabla_\theta H_{\hat\beta_n}(\theta)\right]
\neq
\nabla_\theta H_\beta(\theta).
\]
\paragraph{One-step model using Wasserstein flow of discrepancy.}
Let $\zeta$ be a simple latent distribution and let $\alpha_\theta=(G_\theta)_\sharp\zeta$ be the model distribution. Assume that the target data distribution is $\beta$. The Wasserstein-flow construction chooses a discrepancy

\(\mathcal E_\beta(\alpha),\)
\begin{equation}\label{eq-one-step-wgf}
\partial_t\al_t+\diverg(\al_t w_t)=0,
\qquad
w_t(x)=-\nabla\delta_\alpha \mathcal E_\beta(\al_t)(x).
\end{equation}
\begin{equation}\label{eq-one-step-l2-fit}
\min_\eta \int
\norm{U_\eta(x)-w_t(x)}^2
\,\d\al_t(x).
\end{equation}
\[
\alpha_{\theta}^{+}
=
(\Id+\tau U_{\eta_t})_\sharp \alpha_\theta,
\qquad\text{or equivalently}\qquad
G_\theta^{+}(z)=G_\theta(z)+\tau U_{\eta_t}(G_\theta(z)).
\]
\paragraph{Sliced-Wasserstein flow.}
\(\mathcal E_\beta(\alpha)=\frac12\SW_2^2(\alpha,\beta),\)
with $\SW_2$ defined in Section~\ref{sec-sliced-wasserstein}. With the continuity-equation convention $\partial_t\alpha_t+\diverg(\alpha_t w_t)=0$, the descent velocity points from the current projected law toward the target projected law.

\begin{equation}\label{eq-sliced-wasserstein-flow-velocity}
w_{\SW}[\alpha,\beta](x)
=
\int_{\Sphere^{d-1}}
\big(T_\theta(P_\theta x)-P_\theta x\big)\,\theta\,\d\sigma(\theta).
\end{equation}
The sign follows from the one-dimensional identity: the first variation of $\frac12\Wass_2^2(\alpha,\beta)$ has spatial derivative $s-T(s)$, so the descent velocity is $T(s)-s$, and composing with $P_\theta$ lifts it in the direction $\theta$. Thus empirical implementations only require one-dimensional sorting along sampled directions, followed by averaging the lifted projected displacements.

\paragraph{Stein variational gradient descent.}\label{sec-svgd-generative-flow}
For the formal density-level calculation, assume $\alpha=\rho_\alpha\,\d x$ and $\beta=\rho_\beta\,\d x$ have smooth positive densities, and let $v$ be a smooth compactly supported vector field. For the perturbation $\alpha_\epsilon=(\Id+\epsilon v)_\sharp\alpha$, integration by parts gives the Stein first variation

\[
\left.\frac{\d}{\d\epsilon}\KL(\alpha_\epsilon|\beta)\right|_{\epsilon=0}
=
-\int
\big(\dotp{\nabla\log\rho_\beta(x)}{v(x)}+\diverg v(x)\big)
\d\alpha(x).
\]
\begin{equation}\label{eq-svgd-velocity}
v_{\alpha}^{\mathrm{SVGD}}(x)
=
\int
\Big(k(y,x)\nabla\log\rho_\beta(y)+\nabla_y k(y,x)\Big)
\d\alpha(y).
\end{equation}
\begin{equation}\label{eq-svgd-mean-field}
\partial_t\alpha_t+\diverg\bigl(\alpha_t v_{\alpha_t}^{\mathrm{SVGD}}\bigr)=0.
\end{equation}
For particles $\alpha^\ell_n=n^{-1}\sum_i\delta_{X_i^\ell}$, this becomes the explicit update

\begin{equation}\label{eq-svgd-particle-update}
X_i^{\ell+1}
=
X_i^\ell
+
\tau\,\frac{1}{n}\sum_{j=1}^n
\Big(k(X_j^\ell,X_i^\ell)\nabla\log\rho_\beta(X_j^\ell)
+
\nabla_{X_j}k(X_j^\ell,X_i^\ell)\Big).
\end{equation}
\paragraph{Self-corrected drifting fields.}
\begin{equation}\label{eq-normalized-kernel-drift}
B_\epsilon[\nu](x)
\eqdef
\frac{\int (y-x)K_\epsilon(x,y)\,\d\nu(y)}
{\int K_\epsilon(x,y)\,\d\nu(y)}.
\end{equation}
For the Gaussian kernel $K_\epsilon(x,y)=\exp(-\norm{x-y}^2/(2\epsilon))$, this normalized field is a score of a smoothed density:

\begin{equation}\label{eq-normalized-kernel-score}
B_\epsilon[\nu](x)
=
\epsilon\nabla\log\!\left(\int K_\epsilon(x,y)\,\d\nu(y)\right).
\end{equation}
\begin{equation}\label{eq-cross-minus-self-drift}
u_t(x)=B_\epsilon[\beta](x)-B_\epsilon[\al_t](x)
=
\epsilon\nabla\log
\frac{\int K_\epsilon(x,y)\,\d\beta(y)}
{\int K_\epsilon(x,y)\,\d\al_t(y)}.
\end{equation}
The first term pulls samples toward data, while the second term corrects self-attraction and prevents all particles from collapsing onto the same high-density region. For a fixed reference measure, \(B_\epsilon[\nu]\) is precisely the Gaussian mean-shift displacement in~\eqref{eq-l2-attention-mean-shift}: it moves $x$ toward the local kernel barycenter of $\nu$.

\begin{prop}[Instantaneous gradient representation of drifting]\label{prop-drifting-semi-relaxed-gradient}
Let $\al_t=\rho_t\,\d x$ be a smooth curve of probability measures with positive densities, and let $u_t=\nabla\phi_t$ be a smooth time-dependent gradient field. Define the semi-relaxed functional
\begin{equation}\label{eq-semi-relaxed-drift-functional}
\mathcal R_t(\alpha|\al_t)
\eqdef
-\int \phi_t(x)\,\d\alpha(x)
+\int \phi_t(x)\,\d\al_t(x).
\end{equation}
Here $\al_t$ and $\phi_t$ are frozen when taking the first variation with respect to the first argument $\alpha$. Then the continuity equation
\(\partial_t\al_t+\diverg(\al_t u_t)=0\)
is the formal Wasserstein gradient descent of the frozen time-dependent functional $\alpha\mapsto\mathcal R_t(\alpha|\al_t)$.
\end{prop}
\begin{proof}
Since $\al_t$ and $\phi_t$ are fixed in the variation with respect to $\alpha$, the first variation is \(\delta_\alpha \mathcal R_t(\alpha|\al_t)(x)=-\phi_t(x).\)
By Proposition~\ref{prop-formal-wass-gradient}, \(\Wgrad \mathcal R_t(\alpha|\al_t) = \nabla\delta_\alpha \mathcal R_t(\alpha|\al_t) = -\nabla\phi_t = -u_t.\)
The Wasserstein gradient-descent velocity is the negative of this gradient, namely $u_t$. Substituting this velocity in the continuity equation gives the claimed flow.
\end{proof}
\subsection{Moment Measures}
\label{sec-moment-measures}
\begin{defn}[Moment measure of a convex function]\label{def-moment-measure}
Let $u:\RR^d\to\RR\cup\{+\infty\}$ be a proper lower-semicontinuous convex function such that
\(Z_u\eqdef \int_{\RR^d} e^{-u(x)}\,\d x <+\infty.\)
The normalized log-concave measure associated with $u$ is \(\eta_u \eqdef Z_u^{-1} e^{-u(x)}\,\d x.\)
Whenever $\nabla u$ is defined $\eta_u$-almost everywhere, the moment measure of $u$ is \(\mathfrak M(u)\eqdef (\nabla u)_\sharp \eta_u.\)
\end{defn}
\(\int y\,\d\mathfrak M(u)(y) = Z_u^{-1}\int \nabla u(x)e^{-u(x)}\,\d x = -Z_u^{-1}\int \nabla(e^{-u(x)})\,\d x =0.\)
Thus moment measures are necessarily centered. The nonsmooth theory uses essentially continuous convex functions: these are lower-semicontinuous convex functions whose set of discontinuity points has zero $\mathcal H^{d-1}$ measure.

\begin{thm}[Cordero--Erausquin--Klartag]\label{thm-moment-measure-characterization}
Let $\al\in\Pp_1(\RR^d)$ be a probability measure. There exists an essentially continuous convex function $u$ such that $\al=\mathfrak M(u)$ if and only if
\(\int y\,\d\al(y)=0 \qquad\text{and}\qquad \supp(\al)\text{ is not contained in a hyperplane.}\)
Under these conditions, $u$ is unique up to translations of the argument and additive constants.
\end{thm}
\paragraph{Optimal-transport variational formulation.}
\begin{equation}\label{eq-moment-max-correlation}
\mathcal C_\al(\eta)
\eqdef
\sup_{\pi\in\Couplings(\eta,\al)}
\int_{\RR^d\times\RR^d} x\cdot y\,\d\pi(x,y).
\end{equation}
\begin{equation}\label{eq-moment-max-correlation-dual}
\mathcal C_\al(\eta)
=
\inf_{v\ \mathrm{convex}}
\left\{
\int v(x)\,\d\eta(x)+\int v^*(y)\,\d\al(y)
\right\},
\end{equation}
where the infimum is over convex functions for which both integrals are well defined and $v^*$ is the Legendre transform. If $\eta,\al\in\Pp_2(\RR^d)$, then

\begin{equation}\label{eq-moment-correlation-w2}
\mathcal C_\al(\eta)
=
\frac12\int \norm{x}^2\,\d\eta(x)
+
\frac12\int \norm{y}^2\,\d\al(y)
-
\frac12\Wass_2^2(\eta,\al).
\end{equation}
\begin{equation}\label{eq-moment-measure-variational}
\min_{\eta\in\Pp_1(\RR^d)}
\left\{
\mathcal H(\eta)+\mathcal C_\al(\eta)
\right\},
\qquad
\mathcal H(r\,\d x)\eqdef \int r(x)\log r(x)\,\d x,
\end{equation}
\begin{prop}[Variational characterization of moment measures]\label{prop-moment-hidden-convexity}
Let $\al\in\Pp_1(\RR^d)$ be centered and not supported on a hyperplane. Then the minimization problem~\eqref{eq-moment-measure-variational} admits a solution, unique up to translations. Every minimizer is a probability measure with log-concave density of the form
\(\eta=Z_u^{-1}e^{-u}\,\d x,\)
where $u$ is convex, essentially continuous, and satisfies the moment-measure equation
\(\al=(\nabla u)_\sharp \eta.\)
Conversely, any essentially continuous convex $u$ satisfying this equation yields a global minimizer. If $\al\in\Pp_2(\RR^d)$, the objective $\eta\mapsto\mathcal H(\eta)+\mathcal C_\al(\eta)$ is displacement convex along $\Wass_2$ geodesics of absolutely continuous measures in $\Pp_2(\RR^d)$ whenever the terms are finite. When merely $\al\in\Pp_1(\RR^d)$, the maximal-correlation term remains convex along such finite-second-moment geodesics by approximation; existence and uniqueness on the full $\Pp_1$ domain follow from the lower-semicontinuity argument of Santambrogio.
\end{prop}
\begin{proof}
The existence proof has two ingredients. First, since $\al$ is centered, translating $\eta$ does not change either $\mathcal H(\eta)$ or $\mathcal C_\al(\eta)$, so one can center a minimizing sequence. Second, the assumption that $\al$ is not supported on a hyperplane gives a coercive lower bound of the form $\mathcal C_\al(\eta)\geq c_\al\int\norm{x}\,\d\eta(x)$ for centered absolutely continuous $\eta$, with $c_\al>0$. Together with the lower-semicontinuity estimates for entropy and maximal correlation, this yields a minimizer.

Let $\eta=r\,\d x$ be a minimizer and let $u$ be a convex optimizer in the dual formula~\eqref{eq-moment-max-correlation-dual}. Keeping $u$ fixed and varying $\eta$ in~\eqref{eq-moment-measure-variational} gives the Euler equation
\(\log r(x)+1+u(x)=\mathrm{constant} \qquad\text{on }\{r>0\},\)
so $\eta=Z_u^{-1}e^{-u}\d x$. The optimality condition for the scalar-product transport problem says that an optimal coupling is supported on $\{(x,y):y\in\partial u(x)\}$. Since $\eta$ is absolutely continuous and $u$ is convex, $\partial u(x)=\{\nabla u(x)\}$ for $\eta$-almost every $x$, hence $\al=(\nabla u)_\sharp\eta$.

Conversely, assume $\eta=Z_u^{-1}e^{-u}\d x$ and $\al=(\nabla u)_\sharp\eta$. Let $\nu$ be a smooth compactly supported competitor, and let $T$ be the Brenier map from $\eta$ to $\nu$. Along the geodesic $\eta_t=((1-t)\Id+tT)_\sharp\eta$, the right derivative of the entropy at $t=0$ is
\(\frac{\d}{\d t}\mathcal H(\eta_t)\Big|_{t=0^+} = -\int \dotp{T(x)-x}{\nabla u(x)}\,\d\eta(x),\)
where the identity follows by differentiating the Jacobian formula and integrating by parts. The dual optimizer $u$ gives the upper directional bound
\(\frac{\d}{\d t}\mathcal C_\al(\eta_t)\Big|_{t=0^+} \leq \int \dotp{T(x)-x}{\nabla u(x)}\,\d\eta(x).\)
Conversely, Santambrogio's derivative estimate gives \(\frac{\d}{\d t}\mathcal C_\al(\eta_t)\Big|_{t=0^+} \geq \int \dotp{T(x)-x}{\nabla u(x)}\,\d\eta(x),\)
because $(\Id,\nabla u)_\sharp\eta$ is optimal for the scalar-product problem. The two bounds coincide, so the first-order terms cancel. Hence the one-sided derivative of $\mathcal H+\mathcal C_\al$ at $\eta$ in every such direction is zero. Displacement convexity implies global minimality, and approximation removes the smooth compact-support restriction. Strict displacement convexity of entropy gives uniqueness, except for translations; translations do not change $\mathcal C_\al$ because $\al$ is centered.

The entropy term is displacement convex by McCann's theorem, recalled in Theorem~\ref{thm-mccann-internal-energy}. If $\al$ has finite second moment, identity~\eqref{eq-moment-correlation-w2} writes $\mathcal C_\al$ as the sum of the $1$-convex moment term $\eta\mapsto\frac12\int\norm{x}^2\,\d\eta$ and the $(-1)$-convex term $\eta\mapsto-\frac12\Wass_2^2(\eta,\al)$, hence $\mathcal C_\al$ is displacement convex. For a target with only a finite first moment, Santambrogio obtains the same convexity along $\Pp_2$ geodesics by approximation and proves the full variational characterization by lower semicontinuity.
\end{proof}
\paragraph{Conjugate moment measures for generation.}
\begin{equation}\label{eq-conjugate-moment-measure}
\beta
=
(\nabla w^*)_\sharp
\left(Z_w^{-1}e^{-w(z)}\,\d z\right).
\end{equation}
\subsection{Evolution in Depth of Transformers}
\label{sec-transformer-depth-evolution}
\paragraph{Attention as a context-dependent velocity.}
\begin{equation}\label{eq-transformer-residual-attention-step}
x_i
\longmapsto
x_i+\frac1T\operatorname{Att}^{(n)}_\theta(x_1,\ldots,x_n)_i
=
x_i+\frac1T
\sum_{j=1}^n
\frac{e^{\dotp{Qx_i}{Kx_j}}Vx_j}
{\sum_{\ell=1}^n e^{\dotp{Qx_i}{Kx_\ell}}}.
\end{equation}
\paragraph{Token measure evolution.}
Let \(x_i^{(\ell)}\) be token \(i\) after layer \(\ell\), and let \(\al^{(\ell)}=n^{-1}\sum_i\de_{x_i^{(\ell)}}\) be the corresponding empirical law. With \(\tau=1/T\), equation~\eqref{eq-transformer-residual-attention-step} is exactly

\begin{equation}\label{eq-transformer-residual-measure-step}
\al^{(\ell+1)}
=
\left(
\Id+\tau\Gamma_{\theta^{(\ell)}}[\al^{(\ell)}]
\right)_\sharp
\al^{(\ell)},
\qquad
\ell=0,\ldots,T-1.
\end{equation}
\begin{equation*}
\partial_t\al_t
+
\diverg\!\left(\al_t\Gamma_{\theta_t}[\al_t]\right)
=
0.
\end{equation*}
\paragraph{\texorpdfstring{$L^2$ attention and mean shift.}{L^2 attention and mean shift.}}
\[
K_\epsilon(x,y)=\exp(-\norm{x-y}^2/(2\epsilon)),
\qquad
\rho_\epsilon[\alpha](x)=\int K_\epsilon(x,y)\d\alpha(y),
\qquad
m_\epsilon[\alpha](x)
=
\frac{\int yK_\epsilon(x,y)\d\alpha(y)}
{\rho_\epsilon[\alpha](x)}.
\]
\begin{equation}\label{eq-l2-attention-mean-shift}
M_\epsilon[\alpha](x)
\eqdef
m_\epsilon[\alpha](x)-x
=
\frac{\int (y-x)K_\epsilon(x,y)\d\alpha(y)}
{\rho_\epsilon[\alpha](x)}
=
\epsilon\nabla\log\rho_\epsilon[\alpha](x)
\end{equation}
\[
x_i^{(\ell+1)}
=
(1-\tau)x_i^{(\ell)}+\tau m_\epsilon[\alpha^{(\ell)}](x_i^{(\ell)})
=
x_i^{(\ell)}+\tau M_\epsilon[\alpha^{(\ell)}](x_i^{(\ell)})
\]
\(\partial_t\alpha_t+\diverg\bigl(\alpha_tM_\epsilon[\alpha_t]\bigr)=0.\)
\paragraph{Consensus and Markov averaging.}
\begin{equation}\label{eq-general-mean-shift-field}
M_K[\alpha](x)
\eqdef
\frac{\int (y-x)K(x,y)\d\alpha(y)}
{\int K(x,y)\d\alpha(y)}.
\end{equation}
For an empirical law $\alpha=\sum_{i=1}^n a_i\delta_{x_i}$, where $a_i>0$ and $\sum_i a_i=1$, let $X\in\RR^{n\times d}$ have rows $x_i^\top$, set $\a=(a_i)_i$, and define

\begin{equation}\label{eq-mean-shift-markov-matrix}
(\K_X)_{i,j}=K(x_i,x_j),
\qquad
\mathsf P_X
\eqdef
\diag(\K_X\a)^{-1}\K_X\diag(\a).
\end{equation}
\begin{equation}\label{eq-mean-shift-particle-system}
\dot x_i
=
\frac{\sum_j a_jK(x_i,x_j)(x_j-x_i)}
{\sum_j a_jK(x_i,x_j)},
\qquad\text{or equivalently}\qquad
\dot X=(\mathsf P_X-I_n)X.
\end{equation}
\paragraph{Dobrushin contraction from Hilbert geometry.}
\(\Hilbert(e^z,e^{z'})=\norm{z-z'}_V.\)
A row-stochastic matrix $\mathsf P$ is order preserving and satisfies $\mathsf P(z+s\ones_n)=\mathsf Pz+s\ones_n$. It is therefore a linear topical map, so Proposition~\ref{prop-topical-variation-nonexpansive} gives nonexpansiveness on the additive quotient.

\begin{prop}[Dobrushin contraction is the tangent form of Birkhoff contraction]\label{prop-dobrushin-birkhoff-contraction}
Let $\mathsf P\in\RR_+^{n\times n}$ be row-stochastic. Its Dobrushin coefficient is
\begin{equation}\label{eq-dobrushin-coefficient}
\delta(\mathsf P)
\eqdef
\frac12\max_{i,\ell}\sum_j\abs{\mathsf P_{i,j}-\mathsf P_{\ell,j}}
=
1-\min_{i,\ell}\sum_j\min\{\mathsf P_{i,j},\mathsf P_{\ell,j}\}.
\end{equation}
It is the exact operator norm induced by $\norm{\cdot}_V$ on $\RR^n/\Span(\ones_n)$:
\begin{equation}\label{eq-dobrushin-exact-quotient-norm}
\delta(\mathsf P)
=
\sup_{z\notin\Span(\ones_n)}
\frac{\norm{\mathsf Pz}_V}{\norm{z}_V}.
\end{equation}
\eqllead{Consequently,}{eq-dobrushin-variation-contraction}{
\norm{\mathsf Pz}_V
\leq
\delta(\mathsf P)\norm{z}_V
\qquad (z\in\RR^n).
}
If $\mathsf P>0$, then this exact tangent factor is bounded by the Birkhoff coefficient of Theorem~\ref{thm-birkhoff}:
\begin{equation}\label{eq-dobrushin-birkhoff-comparison}
\delta(\mathsf P)\leq\la(\mathsf P)<1.
\end{equation}
\end{prop}
\begin{proof}
Fix two rows $p=(\mathsf P_{i,j})_j$ and $p'=(\mathsf P_{\ell,j})_j$, and remove their common mass $r_j=\min\{p_j,p'_j\}$. The two residual nonnegative measures have the same mass
\(1-\sum_jr_j = \frac12\sum_j\abs{p_j-p'_j}\).
Their averages of $z$ can differ by at most this mass times $\norm{z}_V$. Maximizing over the two output rows proves the upper bound in~\eqref{eq-dobrushin-exact-quotient-norm}. Conversely, for a pair of rows attaining the maximum, choose $z_j=1$ when $p_j\geq p'_j$ and $z_j=0$ otherwise. Then $\norm{z}_V=1$ and $(\mathsf Pz)_i-(\mathsf Pz)_\ell=\frac12\sum_j\abs{p_j-p'_j}$, proving equality and hence~\eqref{eq-dobrushin-variation-contraction}.

Assume now that $\mathsf P>0$. Apply Theorem~\ref{thm-birkhoff} to $u_s=e^{sz}$ and $\ones_n$. Since $\mathsf P\ones_n=\ones_n$,
\(\Hilbert(\mathsf P e^{sz},\ones_n) \leq \la(\mathsf P)\Hilbert(e^{sz},\ones_n) = \la(\mathsf P)|s|\norm{z}_V.\)
Moreover, $\mathsf P e^{sz}=\ones_n+s\mathsf Pz+O(s^2)$, hence $\log(\mathsf P e^{sz})=s\mathsf Pz+O(s^2)$. Divide by $|s|$ and let $s\to0$ to obtain $\norm{\mathsf Pz}_V\leq\la(\mathsf P)\norm{z}_V$. Taking the exact quotient norm~\eqref{eq-dobrushin-exact-quotient-norm} yields~\eqref{eq-dobrushin-birkhoff-comparison}. This Hilbert--Hopf--Dobrushin relation extends to Markov operators on general cones.
\end{proof}
\(\norm{\mathsf P^\top\p-\mathsf P^\top\p'}_{\ell^1} \leq \delta(\mathsf P)\norm{\p-\p'}_{\ell^1}.\)
Thus~\eqref{eq-dobrushin-variation-contraction} contracts observables modulo constants, while the adjoint inequality contracts probability laws in total variation. The proposition also separates two useful constants.

The same coefficient applies to the measure-valued dynamics. For $\alpha\in\Pp(C)$, where $C\subset\RR^d$ is compact, define the Markov operator

\begin{equation}\label{eq-mean-shift-markov-operator}
(\mathsf P_\alpha h)(x)
\eqdef
\int_C h(y)p_{\alpha,x}(y)\d\alpha(y),
\qquad
p_{\alpha,x}(y)
\eqdef
\frac{K(x,y)}{Z_\alpha(x)},
\qquad
Z_\alpha(x)
\eqdef
\int_C K(x,z)\d\alpha(z).
\end{equation}
\begin{equation}\label{eq-mean-shift-continuous-dobrushin}
\delta(\mathsf P_\alpha)
\eqdef
\frac12\sup_{x,x'\in\supp\alpha}
\int_C\abs{p_{\alpha,x}(y)-p_{\alpha,x'}(y)}\d\alpha(y).
\end{equation}
\(\norm{\mathsf P_\alpha h}_{V,\supp\alpha} \leq \delta(\mathsf P_\alpha)\norm{h}_{V,\supp\alpha}.\)
\begin{equation}\label{eq-mean-shift-kernel-birkhoff-factor}
\eta_K(C)
\eqdef
\sup_{x,x',y,y'\in C}
\frac{K(x,y)K(x',y')}{K(x',y)K(x,y')},
\qquad
\la_K(C)
\eqdef
\frac{\sqrt{\eta_K(C)}-1}{\sqrt{\eta_K(C)}+1}.
\end{equation}
\begin{equation}\label{eq-mean-shift-cross-ratio}
\frac{(\mathsf P_X)_{i,j}(\mathsf P_X)_{\ell,r}}
{(\mathsf P_X)_{i,r}(\mathsf P_X)_{\ell,j}}
=
\frac{K(x_i,x_j)K(x_\ell,x_r)}
{K(x_i,x_r)K(x_\ell,x_j)}.
\end{equation}
\begin{equation}\label{eq-mean-shift-projective-chain}
\delta(\mathsf P_X)
\leq
\la(\mathsf P_X)
=
\la(\K_X)
\leq
\la_K(C).
\end{equation}
\begin{equation}\label{eq-mean-shift-uniform-dobrushin}
\bar\delta_K(C)
\eqdef
\sup_{\alpha\in\Pp(C)}\delta(\mathsf P_\alpha)
\leq
\la_K(C)<1.
\end{equation}
\begin{equation}\label{eq-mean-shift-coarse-birkhoff-bound}
\la_K(C)
\leq
\frac{k_+-k_-}{k_++k_-}.
\end{equation}
\begin{thm}[Dobrushin consensus for positive mean shift]\label{thm-mean-shift-consensus}
Let $\alpha_0\in\Pp(\RR^d)$ have compact support, and set
\(C_0=\operatorname{conv}(\supp\alpha_0), \qquad D_0=\operatorname{diam}(\supp\alpha_0), \qquad \bar\delta=\bar\delta_K(C_0)<1.\)
Assume that $K:C_0\times C_0\to(0,+\infty)$ is Lipschitz.
For an empirical initial measure $\alpha_0=\sum_i a_i\delta_{x_i^{(0)}}$, let $X^{(\ell)}$ follow the damped blurring iteration
\(X^{(\ell+1)} = \bigl((1-\tau)I_n+\tau\mathsf P_{X^{(\ell)}}\bigr)X^{(\ell)}, \qquad 0<\tau\leq1,\)
define
\(\delta^{(\ell)}=\delta(\mathsf P_{X^{(\ell)}}), \qquad q^{(\ell)}=1-\tau(1-\delta^{(\ell)}), \qquad q_\tau=1-\tau(1-\bar\delta)<1.\)
\eqllead{Then}{eq-mean-shift-discrete-diameter-contraction}{
D(X^{(\ell)})
\eqdef
\max_{i,j}\norm{x_i^{(\ell)}-x_j^{(\ell)}}
\leq
D_0\prod_{r=0}^{\ell-1}q^{(r)}
\leq
q_\tau^\ell D_0.
}
All particles converge to a common $x_\infty^{\rm d}\in C_0$, and
\(\max_i\norm{x_i^{(\ell)}-x_\infty^{\rm d}} \leq \frac{D_0}{1-\bar\delta}q_\tau^\ell.\)
The characteristic solution of
\begin{equation}\label{eq-positive-mean-shift-pde}
\partial_t\alpha_t+\diverg\bigl(\alpha_tM_K[\alpha_t]\bigr)=0,
\qquad
\alpha_{t=0}=\alpha_0,
\end{equation}
is globally defined and has nested convex hulls, \(\operatorname{conv}(\supp\alpha_t) \subset \operatorname{conv}(\supp\alpha_s) \qquad (t\geq s\geq0).\)
Writing $D(t)=\operatorname{diam}(\supp\alpha_t)$, one has the adaptive estimate
\begin{equation}\label{eq-mean-shift-diameter-contraction}
D(t)
\leq
D(s)\exp\!\left(-\int_s^t\bigl[1-\delta(\mathsf P_{\alpha_r})\bigr]\d r\right)
\leq
e^{-(1-\bar\delta)(t-s)}D(s)
\qquad (t\geq s\geq0).
\end{equation}
\begin{equation}\label{eq-mean-shift-dirac-convergence}
\Wass_\infty(\alpha_t,\delta_{x_\infty^{\rm c}})
\leq
\frac{D_0}{1-\bar\delta}e^{-(1-\bar\delta)t}.
\end{equation}
\end{thm}
\begin{proof}
Because $K$ is positive and Lipschitz on the compact set $C_0\times C_0$, its minimum is positive. Hence the denominator in~\eqref{eq-general-mean-shift-field} is uniformly bounded away from zero, and the characteristic vector field is uniformly Lipschitz in $x$ and Lipschitz in $\alpha$ for $\Wass_1$. Moreover, $x+M_K[\alpha](x)=\mathsf P_\alpha\Id(x)$ lies in $\operatorname{conv}(\supp\alpha)$. No characteristic can cross an outward supporting hyperplane of the current convex hull, so the convex hulls are nested and the dynamics is globally well posed.

For any compact set $S\subset\RR^d$, the support-function identity
\begin{equation}\label{eq-mean-shift-directional-diameter}
\operatorname{diam}(S)
=
\sup_{\theta\in\Sphere^{d-1}}
\left(\sup_{x\in S}\dotp{\theta}{x}-\inf_{x\in S}\dotp{\theta}{x}\right)
\end{equation}
reduces Euclidean diameter contraction to scalar variation contraction. For a finite configuration, set $z_\theta=X\theta\in\RR^n$. Proposition~\ref{prop-dobrushin-birkhoff-contraction} gives
\begin{align*}
\norm{z_\theta^{(\ell+1)}}_V
&\leq
(1-\tau)\norm{z_\theta^{(\ell)}}_V
+\tau\norm{\mathsf P_{X^{(\ell)}}z_\theta^{(\ell)}}_V\\
&\leq
\bigl[1-\tau(1-\delta(\mathsf P_{X^{(\ell)}}))\bigr]\norm{z_\theta^{(\ell)}}_V
\leq
q_\tau\norm{z_\theta^{(\ell)}}_V.
\end{align*}
Taking the supremum over $\theta$ gives the one-step factor $q^{(\ell)}$; multiplying these factors and using $q^{(\ell)}\leq q_\tau$ proves~\eqref{eq-mean-shift-discrete-diameter-contraction}. Since $\norm{x_i^{(\ell+1)}-x_i^{(\ell)}}\leq\tau D(X^{(\ell)})$, every particle path is Cauchy; the vanishing diameter forces a common limit, and summing the uniform geometric tail yields the displayed pointwise estimate.

For the measure-valued flow, let $f_\theta(x)=\dotp{\theta}{x}$ and denote by $w_\theta(t)=\norm{f_\theta}_{V,\supp\alpha_t}$ the width of the support in direction $\theta$. The upper Dini derivative of the extrema along characteristics satisfies
\[
D^+w_\theta(t)
\leq
\norm{\mathsf P_{\alpha_t}f_\theta}_{V,\supp\alpha_t}
-w_\theta(t)
\leq
-\bigl[1-\delta(\mathsf P_{\alpha_t})\bigr]w_\theta(t).
\]
Gronwall's lemma, followed by the support-function identity~\eqref{eq-mean-shift-directional-diameter}, proves~\eqref{eq-mean-shift-diameter-contraction}. Finally, every characteristic satisfies $\norm{\dot X_t}\leq D(t)$. It is therefore Cauchy, all characteristic limits coincide because $D(t)\to0$, and integrating the uniform exponential bound from $t$ to $+\infty$ gives~\eqref{eq-mean-shift-dirac-convergence}.
\end{proof}
For the Gaussian kernel $K_\epsilon(x,y)=e^{-\norm{x-y}^2/(2\epsilon)}$, the cross-ratio has the closed form

\(\log\frac{K_\epsilon(x,y)K_\epsilon(x',y')} {K_\epsilon(x',y)K_\epsilon(x,y')} = \frac{\dotp{x-x'}{y-y'}}{\epsilon}.\)
Hence, for a compact convex set $C$ of diameter $D$,

\begin{equation}\label{eq-gaussian-mean-shift-birkhoff-factor}
\eta_{K_\epsilon}(C)=e^{D^2/\epsilon},
\qquad
\la_{K_\epsilon}(C)=\tanh\!\left(\frac{D^2}{4\epsilon}\right).
\end{equation}
\[
D(X^{(\ell+1)})
\leq
\tanh\!\left(\frac{D(X^{(\ell)})^2}{4\epsilon}\right)D(X^{(\ell)})
\leq
\frac{D(X^{(\ell)})^3}{4\epsilon},
\]
\begin{equation}\label{eq-gaussian-mean-shift-adaptive-contraction}
D(t)
\leq
D(s)\exp\!\left(
-\int_s^t
\left[1-\tanh\!\left(\frac{D(r)^2}{4\epsilon}\right)\right]\d r
\right).
\end{equation}
\paragraph{Scope of the consensus result.}
The projective argument also clarifies what is not conserved and when clustering replaces consensus. Because row normalization makes the pairwise coefficients asymmetric, the consensus point need not equal the barycenter of $\alpha_0$; that barycenter is conserved for the constant kernel, but not in general.

\paragraph{Gradient structure and limitations.}
When the token space has dimension $d$ and the query/key space has dimension $r$, take $Q,K\in\RR^{r\times d}$ and $V\in\RR^{d\times d}$. If $V=Q^\top K$, the field $\Gamma_\theta[\alpha]$ is a gradient vector field in the token variable. Indeed, define the log-partition potential

\(\Phi_\alpha(x) = \int \exp(\dotp{Qx}{Ky})\d\alpha(y), \qquad U_\alpha(x)=\log\Phi_\alpha(x).\)
\(\nabla_x U_\alpha(x) = \frac{\int Q^\top K y\,\exp(\dotp{Qx}{Ky})\d\alpha(y)} {\int \exp(\dotp{Qx}{Kz})\d\alpha(z)} = \Gamma_\theta[\alpha](x).\)
\(f_{\rm att}(\alpha)=\int U_\alpha(x)\,\d\alpha(x)\)
\(\delta f_{\rm att}(\alpha)(z) = U_\alpha(z) + \int \frac{\exp(\dotp{Qx}{Kz})}{\Phi_\alpha(x)} \,\d\alpha(x),\)
up to an additive constant. The second term is the response of every query normalization to a perturbation of the key distribution, and its spatial gradient is absent from $\Gamma_\theta[\alpha]$.

\subsection{Flows over the Gaussian Manifold}
\label{sec-gaussian-closure-transport-dynamics}
\paragraph{Gaussianity preservation.}
\begin{prop}[Affine velocities preserve Gaussianity]\label{prop-gaussian-affine-closure}
Let $\alpha_0=\Gaussian(m_0,\Sigma_0)$, with $\Sigma_0\succ0$. Let $b_t\in\RR^d$ and $A_t\in\RR^{d\times d}$ be locally integrable on a time interval, and let $(m_t,\Sigma_t)$ solve
\(\dot m_t=b_t, \qquad \dot\Sigma_t=A_t\Sigma_t+\Sigma_tA_t^\top, \qquad (m_{t=0},\Sigma_{t=0})=(m_0,\Sigma_0).\)
Then, as long as this matrix ODE is defined, $\Sigma_t\succ0$ and
\(\alpha_t=\Gaussian(m_t,\Sigma_t)\)
is the solution of the continuity equation
\(\partial_t\alpha_t+\diverg(\alpha_t v_t)=0, \qquad v_t(x)=b_t+A_t(x-m_t).\)
In particular, if a smooth functional $f$ has a Wasserstein gradient on Gaussian measures of the affine form
\(\Wgrad f(\Gaussian(m,\Sigma))(x) = b_f(m,\Sigma)+A_f(m,\Sigma)(x-m),\)
with $A_f(m,\Sigma)$ symmetric, then the Wasserstein gradient flow of $f$, initialized from any non-degenerate Gaussian, stays Gaussian and satisfies \(\dot m_t=h_f(m_t,\Sigma_t), \qquad \dot\Sigma_t=H_f(m_t,\Sigma_t),\)
where \(h_f(m,\Sigma)=-b_f(m,\Sigma), \qquad H_f(m,\Sigma) = -\bigl(A_f(m,\Sigma)\Sigma+\Sigma A_f(m,\Sigma)\bigr).\)
Conversely, fix a non-degenerate Gaussian \(\Gaussian(m,\Sigma)\). Suppose that a finite-energy Wasserstein tangent field \(V_{m,\Sigma}\) is represented by a gradient in \(L^2(\Gaussian(m,\Sigma))\) and is tangent to the Gaussian manifold at this Gaussian. Then there exist \(b(m,\Sigma)\) and a symmetric matrix \(A(m,\Sigma)\) such that
\(V_{m,\Sigma}(x)=b(m,\Sigma)+A(m,\Sigma)(x-m) \qquad \Gaussian(m,\Sigma)\text{-a.e.}\)
Without the Wasserstein gradient representative, this converse is false because one may add velocity fields with zero \(\Gaussian(m,\Sigma)\)-weighted divergence.
Finally, any smooth Gaussian curve with positive definite covariance can be generated by an affine velocity. If one wants the velocity to be a Wasserstein tangent gradient, one chooses the unique symmetric solution of the Lyapunov equation
\(A_t\Sigma_t+\Sigma_t A_t=\dot\Sigma_t.\)
\end{prop}
\begin{proof}
Let \(X_t\) follow the characteristic ODE \(\dot X_t=b_t+A_t(X_t-m_t)\) with \(X_0\sim\Gaussian(m_0,\Sigma_0)\). Since \(\dot m_t=b_t\), the centered variable \(\tilde X_t=X_t-m_t\) solves the homogeneous linear ODE \(\dot{\tilde X}_t=A_t\tilde X_t\). If \(\Phi_t\) is the fundamental matrix \(\dot\Phi_t=A_t\Phi_t\), \(\Phi_{t=0}=\Id\), then
\(X_t=m_t+\Phi_t(X_0-m_0), \qquad \Sigma_t=\Phi_t\Sigma_0\Phi_t^\top.\)
Hence \(X_t\) is Gaussian and \(\Sigma_t\succ0\), and
\(\dot\Sigma_t = \frac{\d}{\d t}\EE\bigl(\tilde X_t\tilde X_t^\top\bigr) = A_t\Sigma_t+\Sigma_t A_t^\top.\)
This proves Gaussian preservation and the moment ODE. The Wasserstein gradient-flow statement follows by inserting the descent velocity $v_t=-\Wgrad f(\alpha_t)$.

For the converse, fix $\alpha=\Gaussian(m,\Sigma)$ and denote its density by $\rho$. Tangency to the Gaussian manifold means that the density variation $-\diverg(\rho V)$ is generated by some moment variation $(\dot m,\dot\Sigma)$, with $\dot\Sigma$ symmetric. Set $b=\dot m$, and let $A=A^\top$ be the unique solution of
\(A\Sigma+\Sigma A=\dot\Sigma\).
By the first part of the proposition, the affine gradient field $V_0(x)=b+A(x-m)$ generates exactly the same infinitesimal Gaussian variation. Hence
\(\diverg\bigl(\rho(V-V_0)\bigr)=0\)
in the distributional sense. Both $V$ and $V_0$ belong to the $L^2(\alpha)$ closure of gradient fields. The weighted Helmholtz decomposition therefore makes $V-V_0$ simultaneously a tangent gradient and orthogonal to every tangent gradient, so $V=V_0$ in $L^2(\alpha)$. Equivalently, for a smooth representative $V-V_0=\nabla\psi$, integration by parts gives $\int\norm{\nabla\psi}^2\d\alpha=0$. This proves that the Wasserstein tangent representative is affine. The qualification is essential: without selecting the gradient representative, one can add a nonzero field with zero $\alpha$-weighted divergence without changing the Gaussian curve.

For a prescribed smooth Gaussian curve, set $b_t=\dot m_t$ and choose any matrix $A_t$ satisfying $A_t\Sigma_t+\Sigma_tA_t^\top=\dot\Sigma_t$. Since $\Sigma_t$ is positive definite, the Lyapunov map $A\mapsto A\Sigma_t+\Sigma_t A$ is invertible on symmetric matrices, which gives the unique symmetric choice when a gradient velocity is required. In that case $v_t$ is the gradient of the quadratic potential $x\mapsto \dotp{b_t}{x}+\dotp{A_t(x-m_t)}{x-m_t}/2$.
\end{proof}
\paragraph{Gaussian-preserving gradient flows.}
\begin{prop}[Gaussian closure catalogue]\label{prop-centered-gaussian-covariance-catalogue}
Let \(\gamma=\Gaussian(\bar m,\bar\Sigma)\) on \(\RR^d\), with \(\bar\Sigma\succ0\), and let the initial condition be \(\alpha_0=\Gaussian(m_0,\Sigma_0)\), with \(\Sigma_0\succ0\). For each functional displayed below, let \(\alpha_t\) be the usual Wasserstein gradient flow on \(\Pp_2(\RR^d)\), initialized at \(\alpha_0\). Then the Gaussian family is invariant for these dynamics: as long as the solution exists and \(\Sigma_t\succ0\),
\(\alpha_t=\Gaussian(m_t,\Sigma_t),\)
and the mean and covariance satisfy
\(\dot{m}_t=h(m_t,\Sigma_t), \qquad \dot{\Sigma}_t=H(m_t,\Sigma_t).\)
Write \(\delta_m=m-\bar m\), \(A=\bar\Sigma^{-1}\), and
\(M_{\Sigma,\bar\Sigma} \eqdef \Sigma^{-1/2}\bigl(\Sigma^{1/2}\bar\Sigma\Sigma^{1/2}\bigr)^{1/2}\Sigma^{-1/2}.\)
With the normalizations displayed in the first column, the mean vector field \(h\) and covariance vector field \(H\) are listed in the following table. Gradients with respect to \(\Sigma\) are symmetric Frobenius gradients on the cone of covariance matrices.
\begin{center}
\begingroup
\small
\renewcommand{\arraystretch}{1.55}
\textbf{Gaussian-preserving Wasserstein gradient flows.}\par\smallskip
\endgroup
\end{center}
Here, in the MMD row, \(R=\Sigma+m\,m^\top-\bar\Sigma-\bar m\,\bar m^\top.\)
The debiased Sinkhorn row uses the notation of Corollary~\ref{cor-gaussian-sinkhorn-divergence}: for Gaussian inputs,
\(\bar\MK_{\norm{\cdot-\cdot}^2}^{\epsilon}(\alpha,\gamma) = \norm{\delta_m}^2+\Bb_\epsilon(\Sigma,\bar\Sigma)^2,\)
with the closed-form covariance gradient
\[
G_\epsilon(\Sigma,\bar\Sigma)
=
\tau_\epsilon(\Sigma)
-
\bar\Sigma^{1/2}
\tau_\epsilon\bigl(B_{\Sigma,\bar\Sigma}^{1/2}\bigr)
B_{\Sigma,\bar\Sigma}^{-1/2}
\bar\Sigma^{1/2},
\qquad
B_{\Sigma,\bar\Sigma}=\bar\Sigma^{1/2}\Sigma\bar\Sigma^{1/2}.
\]
Here \(\tau_\epsilon\) is the scalar function
\(\tau_\epsilon(r) \eqdef \frac{\sqrt{\epsilon^2+16r^2}-\epsilon}{4r}, \qquad r>0,\)
applied to positive matrices by spectral calculus. Equivalently, for \(M\succ0\), \(\tau_\epsilon(M) = \bigl(\sqrt{\epsilon^2 \Id+16M^2}-\epsilon \Id\bigr)(4M)^{-1}.\)
With this convention, \(G_\epsilon=\nabla_\Sigma \Bb_\epsilon(\Sigma,\bar\Sigma)^2\).
In the sliced row, \(\sigma\) is the normalized spherical measure on \(\Sphere^{d-1}\), and
\[
G_{\mathrm{sw}}(\Sigma,\bar\Sigma)
=
\int_{\Sphere^{d-1}}
\left(
1-\sqrt{\frac{\theta^\top\bar\Sigma\theta}{\theta^\top\Sigma\theta}}
\right)
\theta\theta^\top\,\d\sigma(\theta).
\]
Here
\[
\mathcal I(\alpha|\gamma)
=
\int \left|\nabla\log\frac{\rho(x)}{\rho_\gamma(x)}\right|^2\rho(x)\,\d x
=
\int
\left|\nabla\log\rho(x)+A(x-\bar m)\right|^2\rho(x)\,\d x
\qquad(\alpha=\rho\,\d x),
\]
where $\rho_\gamma$ is the density of $\gamma$.
\end{prop}
\begin{proof}
For the moment-functional row, the formula is exact on the full Wasserstein space. Indeed, write
\(m_\alpha=\int x\,\d\alpha(x), \qquad \Sigma_\alpha=\int (x-m_\alpha)(x-m_\alpha)^\top\d\alpha(x).\)
If \(\eta\) is a signed perturbation with zero total mass, then
\(\d m_\alpha[\eta]=\int x\,\d\eta(x), \qquad \d\Sigma_\alpha[\eta]=\int (x-m_\alpha)(x-m_\alpha)^\top\d\eta(x).\)
Thus, up to an irrelevant additive constant, the first variation is \(\delta f(\alpha)(x) = \dotp{\nabla_m g}{x} + \dotp{\nabla_\Sigma g}{(x-m_\alpha)(x-m_\alpha)^\top},\)
and its spatial gradient is the affine field
\(\nabla_x\delta f(\alpha)(x) = \nabla_m g+2\nabla_\Sigma g\,(x-m_\alpha).\)

For the quadratic potential row, the ambient first variation is \(\delta f(\alpha)(x)=\frac12x^\top Bx+\dotp{\ell}{x}, \qquad \nabla_x\delta f(\alpha)(x)=B m+\ell+B(x-m),\)
which gives the displayed affine velocity. For the quadratic interaction row, the ambient first variation is \(\delta f(\alpha)(x) = \frac12\int (x-y)^\top G(x-y)\d\alpha(y), \qquad \nabla_x\delta f(\alpha)(x)=G(x-m_\alpha).\)
These first variations are quadratic in \(x\), hence the affine Gaussian flow coincides with the full Wasserstein gradient flow.

For the KL row, write $\alpha=\rho\,\d x$ and let $\rho_\gamma$ denote the density of $\gamma$. The ambient first variation is \(\log(\rho/\rho_\gamma)\) up to an additive constant. If \(\alpha=\Gaussian(m,\Sigma)\), then
\(\nabla\log(\rho/\rho_\gamma)(x) = -\Sigma^{-1}(x-m)+A(x-\bar m),\)
so the descent velocity is
\(v(x)=-A\delta_m+(\Sigma^{-1}-A)(x-m)\).
Proposition~\ref{prop-gaussian-affine-closure} gives \(h(m,\Sigma)=-A\delta_m\) and \(H(m,\Sigma)=2\Id-\Sigma A-A\Sigma\). For the Fisher row, writing \(u=\log(\rho/\rho_\gamma)\), the ambient first variation of \(\frac12\mathcal I(\alpha|\gamma)\) is, up to constants,
\(-\Delta u-\frac12|\nabla u|^2-\dotp{\nabla u}{\nabla\log\rho_\gamma}.\)
For Gaussian \(\rho\) and Gaussian \(\gamma\), this is quadratic in \(x\). Multiplying by two for \(\mathcal I(\alpha|\gamma)\) gives the descent velocity
\(v(x)=-2A^2\delta_m+2(\Sigma^{-2}-A^2)(x-m).\)
Hence the Wasserstein velocity is affine and the Gaussian family is closed, with the displayed \(h\) and \(H\). Although the ambient Fisher flow is a fourth-order PDE for a generic density, it is therefore an exact finite-dimensional closure in this quadratic-reference case.

For \(\Wass_2^2(\alpha,\gamma)\), the Brenier map from \(\Gaussian(m,\Sigma)\) to \(\gamma\) is
\(T(x)=\bar m+M_{\Sigma,\bar\Sigma}(x-m)\).
The descent velocity for the unhalved squared distance is \(2(T-\Id)\), which gives the mean and covariance equations through Proposition~\ref{prop-gaussian-affine-closure}. For the MMD row, \(k(x,y)=\dotp{x}{y}^2\) identifies the kernel mean embedding with the raw second moment. Thus
\(\MMD_k^2(\alpha,\gamma) = \norm{\Sigma+m\,m^\top-\bar\Sigma-\bar m\,\bar m^\top}_{\mathrm F}^2 = \norm{R}_{\mathrm F}^2,\)
and the ambient first variation is \(2x^\top R x\), whose descent velocity is \(-4Rx\).

For the Sinkhorn row, Corollary~\ref{cor-gaussian-sinkhorn-divergence} gives the closed Gaussian formula as a squared mean displacement plus the smoothed Bures term \(\Bb_\epsilon^2\). This is not only a formula on the Gaussian submanifold. The first variation of the entropic value with respect to the first marginal is a Sinkhorn dual potential, and for Gaussian marginals the cross and self dual potentials are quadratic. Their debiased combination is therefore quadratic, so its spatial gradient is affine. Since the restriction of the functional to Gaussians is
\((m,\Sigma)\mapsto \norm{m-\bar m}^2+\Bb_\epsilon(\Sigma,\bar\Sigma)^2,\)
the quadratic coefficient is \(G_\epsilon=\nabla_\Sigma\Bb_\epsilon(\Sigma,\bar\Sigma)^2\), and the moment-functional computation gives the displayed \(h\) and \(H\). To compute this gradient, set \(B_{\Sigma,\bar\Sigma}=\bar\Sigma^{1/2}\Sigma\bar\Sigma^{1/2}\). Since \(\psi_\epsilon'(r)=-2\tau_\epsilon(r)\), differentiating
\(\tr\psi_\epsilon(B_{\Sigma,\bar\Sigma}^{1/2})\) gives
\(-\bar\Sigma^{1/2} \tau_\epsilon\bigl(B_{\Sigma,\bar\Sigma}^{1/2}\bigr) B_{\Sigma,\bar\Sigma}^{-1/2} \bar\Sigma^{1/2},\)
whereas differentiating the self term
\(-\frac12\tr\psi_\epsilon(\Sigma)\) gives \(\tau_\epsilon(\Sigma)\). This proves the displayed closed form for \(G_\epsilon\). When \(\bar\Sigma=\Id\), this covariance contribution is a spectral function of \(\Sigma\). If \(\lambda\) is an eigenvalue of \(\Sigma\), the corresponding scalar velocity is
\(4\sqrt{\lambda+\epsilon^2/16}-4\sqrt{\lambda^2+\epsilon^2/16},\)
which gives the displayed covariance-eigenvalue equation by functional calculus.

For sliced Wasserstein, the exact ambient Wasserstein gradient is obtained by averaging the one-dimensional gradients. Each projection is Gaussian:
\[
(P_\theta)_\sharp\alpha=\Gaussian(\dotp{\theta}{m},\theta^\top\Sigma\theta),
\qquad
(P_\theta)_\sharp\gamma=\Gaussian(\dotp{\theta}{\bar m},\theta^\top\bar\Sigma\theta).
\]
Hence
\[
\SW_2^2(\alpha,\gamma)
=
\int_{\Sphere^{d-1}}
\left[
\dotp{\theta}{\delta_m}^2
+
\left(\sqrt{\theta^\top\Sigma\theta}
-\sqrt{\theta^\top\bar\Sigma\theta}\right)^2
\right]\d\sigma(\theta).
\]
If \(T_\theta\) is the monotone transport from \((P_\theta)_\sharp\alpha\) to \((P_\theta)_\sharp\gamma\), then the descent velocity for the unhalved sliced objective is
\(2\int_{\Sphere^{d-1}}\bigl(T_\theta(P_\theta x)-P_\theta x\bigr)\theta\,\d\sigma(\theta).\)
For Gaussian marginals, \(T_\theta\) is affine. Thus this velocity is affine in \(x\), and Proposition~\ref{prop-gaussian-affine-closure} applies. The spherical identity \(\int\theta\theta^\top\d\sigma(\theta)=\Id/d\) gives the mean equation, and differentiating the covariance term gives \(G_{\mathrm{sw}}\).
\end{proof}
\paragraph{One-dimensional case.}
Fix a Gaussian target $\gamma=\Gaussian(\bar m,\bar\sigma^2)$, with $\bar\sigma>0$, and write $\alpha=\Gaussian(m,\sigma^2)$, with $\sigma>0$. By Proposition~\ref{prop-gaussian-w2-bures}, the parametrization $(m,\sigma)\mapsto\Gaussian(m,\sigma^2)$ identifies the one-dimensional Gaussian family with the open Euclidean half-plane.

\(g_f(m,\sigma)\eqdef f\bigl(\Gaussian(m,\sigma^2)\bigr).\)
Set $\delta=m-\bar m$. For the Sinkhorn row below, with $\epsilon>0$, abbreviate

\[
\psi_\epsilon(r)
\eqdef
-\frac{\sqrt{\epsilon^2+16r^2}-\epsilon}{2}
-\frac{\epsilon}{2}
\log\!\left(
\frac{2\epsilon}{\sqrt{\epsilon^2+16r^2}+\epsilon}
\right),
\qquad r>0.
\]
\paragraph{Constrained evolution on the Gaussian manifold.}
\(\mathcal G=\{\Gaussian(m,\Sigma):m\in\RR^d,\ \Sigma\succ0\}\)
be the Gaussian submanifold of $\Pp_2(\RR^d)$. The Wasserstein gradient of a functional constrained to a smooth submanifold $\mathcal M\subset\Pp_2$ is defined as the Riesz representative of the differential restricted to tangent velocities of $\mathcal M$.

\(\alpha^{(\ell+1)}\in \uargmin{\alpha\in\mathcal M} \frac{1}{2\tau}\Wass_2^2(\alpha,\alpha^{(\ell)})+f(\alpha).\)
For $\mathcal M=\mathcal G$, tangent velocities are affine gradient fields $v(x)=b+A(x-m)$ with $A=A^\top$. The constrained gradient is therefore the $L^2(\Gaussian(m,\Sigma))$ projection of the ambient Wasserstein gradient onto this finite-dimensional affine space, whenever the ambient gradient exists.

\begin{prop}[Gaussian-constrained Wasserstein gradients]\label{prop-gaussian-gradient-bullet-list}
Let $f$ be a smooth functional and assume that its restriction to nondegenerate Gaussian measures can be written as \(f(\Gaussian(m,\Sigma))=F(m,\Sigma).\)
Then the Wasserstein gradient constrained to the Gaussian family is the affine vector field
\(v_F(x) = \nabla_m F(m,\Sigma) + 2\nabla_\Sigma F(m,\Sigma)(x-m),\)
where $\nabla_\Sigma F$ denotes the symmetric matrix derivative. Equivalently, $v_F$ is the $L^2(\Gaussian(m,\Sigma))$ projection of the ambient Wasserstein gradient onto affine gradient fields, whenever the ambient gradient exists. Hence the gradient descent flow constrained to Gaussian measures satisfies
\begin{equation}\label{eq-gaussian-wgf-closure}
\dot m_t=-\nabla_m F(m_t,\Sigma_t),
\qquad
\dot\Sigma_t=-2\bigl(\Sigma_t\nabla_\Sigma F(m_t,\Sigma_t)+\nabla_\Sigma F(m_t,\Sigma_t)\Sigma_t\bigr),
\end{equation}
and the descent velocity is affine.
\end{prop}
\begin{proof}
Test the functional along a Gaussian tangent vector, represented by an affine gradient field
\(v(x)=b+A(x-m)\)
with $A$ symmetric. The induced first-order variations are $\dot m=b$ and $\dot\Sigma=A\Sigma+\Sigma A$. Therefore
\[
\d F(m,\Sigma)[b,A\Sigma+\Sigma A]
=
\dotp{\nabla_m F}{b}
+
\mathrm{tr}\!\left(\nabla_\Sigma F(A\Sigma+\Sigma A)\right).
\]
Since $A$, $\Sigma$ and $\nabla_\Sigma F$ are symmetric, the second term equals
\(2\,\mathrm{tr}(\nabla_\Sigma F\,A\Sigma) = \int \dotp{2\nabla_\Sigma F(x-m)}{A(x-m)}\d\Gaussian(m,\Sigma)(x).\)
Together with the mean term, this gives \(\d F(m,\Sigma)[\dot m,\dot\Sigma] = \int \dotp{v_F(x)}{v(x)}\d\Gaussian(m,\Sigma)(x)\)
for all affine gradient fields $v$. This identifies the constrained Wasserstein gradient in the induced $L^2(\alpha)$ metric, or equivalently the projection of the ambient gradient when it exists. Substituting the descent velocity $-v_F$ in Proposition~\ref{prop-gaussian-affine-closure} gives~\eqref{eq-gaussian-wgf-closure}.
\end{proof}
\paragraph{Non-variational Gaussian-preserving flows.}
\paragraph{Contractive Gaussian projection.}
\begin{defn}[Gaussian projection]\label{def-gaussian-projection}
For \(\alpha\in\Pp_2(\RR^d)\), its Gaussian projection is \(\GaussProj(\alpha) \eqdef \Gaussian(m_\alpha,\Sigma_\alpha),\)
the Gaussian with the same mean and covariance as \(\alpha\).
\end{defn}
\begin{thm}[Gelbrich theorem]\label{thm-gelbrich-projection}
For every $\al,\be\in\Pp_2(\RR^d)$, \(\Wass_2^2\bigl(\GaussProj(\al),\GaussProj(\be)\bigr) = \norm{m_\al-m_\be}^2+\Bb^2(\Sigma_\al,\Sigma_\be) \leq \Wass_2^2(\al,\be).\)
\end{thm}
\begin{proof}
Take any coupling $(X,Y)$ of $\al$ and $\be$, center the variables, and write $C=\EE[(X-m_\al)(Y-m_\be)^\top]$. In the positive definite case, positivity of the block covariance matrix implies the factorization $C=\Sigma_\al^{1/2}K\Sigma_\be^{1/2}$ with $\norm{K}_{\mathrm{op}}\leq1$, and therefore, by operator/nuclear norm duality,
\[
\tr C\leq \tr\left((\Sigma_\al^{1/2}\Sigma_\be\Sigma_\al^{1/2})^{1/2}\right).
\]
The semidefinite case follows by adding $\eta\Id$ to both covariance matrices and letting $\eta\downarrow0$.
Expanding $\EE\norm{X-Y}^2$ gives the lower bound
\(\EE\norm{X-Y}^2 \geq \norm{m_\al-m_\be}^2+\Bb^2(\Sigma_\al,\Sigma_\be).\)
Taking the infimum over couplings proves the inequality, while equality for Gaussian laws is Proposition~\ref{prop-gaussian-w2-bures}.
\end{proof}
\begin{thm}[Hugo Lavenant Gaussian-preservation criterion]\label{thm-lavenant-gaussian-preserving-jko}
Let $f:\Pp_2(\RR^d)\to(-\infty,+\infty]$ satisfy
\(f\bigl(\GaussProj(\al)\bigr)\leq f(\al) \qquad\forall\al\in\Pp_2(\RR^d),\)
with $\GaussProj$ defined in Theorem~\ref{thm-gelbrich-projection}. If $\gamma$ is Gaussian and $\alpha^\star$ minimizes the JKO step
\(\eta\mapsto f(\eta)+\frac1{2\tau}\Wass_2^2(\gamma,\eta),\)
then $\GaussProj(\alpha^\star)$ is also a minimizer. If this JKO minimizer is unique, it is Gaussian.
\end{thm}
\begin{proof}
For the JKO claim, $\GaussProj(\gamma)=\gamma$ because $\gamma$ is Gaussian. Hence, for any competitor $\eta$,
\(f\bigl(\GaussProj(\eta)\bigr)+\frac1{2\tau}\Wass_2^2\bigl(\gamma,\GaussProj(\eta)\bigr) \leq f(\eta)+\frac1{2\tau}\Wass_2^2(\gamma,\eta).\)
Applying this to a minimizer $\eta=\alpha^\star$ shows that $\GaussProj(\alpha^\star)$ is again a minimizer. Uniqueness forces $\alpha^\star=\GaussProj(\alpha^\star)$.
\end{proof}
\paragraph{Moment closure beyond Gaussianity.}
\begin{equation}\label{eq-distribution-free-mean-covariance-closure}
\dot m_t=-\nabla_m g(m_t,\Sigma_t),
\qquad
\dot\Sigma_t
=
-2\bigl(\Sigma_t\nabla_\Sigma g(m_t,\Sigma_t)
+\nabla_\Sigma g(m_t,\Sigma_t)\Sigma_t\bigr).
\end{equation}
For Gaussian initial data, this affine velocity also preserves Gaussianity; for general initial data, the law need not become Gaussian, but its mean and covariance still follow the same closed vector field $(h,H)$ listed in that proposition.

\begin{prop}[Scalar moment closure and eikonal characterization]\label{prop-scalar-moment-closure}
Let $\Omega\subset\RR^d$ be connected and open, let $\varphi\in C^2(\Omega)$, and set
\(m_\varphi(\alpha)\eqdef\int_\Omega\varphi(x)\d\alpha(x)\)
for probability measures such that $\varphi$ and $\norm{\nabla\varphi}^2$ are integrable. For every smooth $g:\RR\to\RR$, the Wasserstein gradient flow of $f_g(\alpha)=g(m_\varphi(\alpha))$ satisfies, whenever the differentiation is justified,
\(\dot m_t = -g'(m_t)\int_\Omega\norm{\nabla\varphi(x)}^2\d\alpha_t(x), \qquad m_t=m_\varphi(\alpha_t).\)
This equation is autonomous for every $g$, meaning that its right-hand side depends on $\alpha_t$ only through $m_t$, if and only if there are constants $a,b\in\RR$ such that
\begin{equation}\label{eq-scalar-moment-affine-carre}
\norm{\nabla\varphi(x)}^2=a+b\varphi(x)
\qquad (x\in\Omega).
\end{equation}
In that case, $\dot m_t=-(a+b m_t)g'(m_t)$.
Moreover, on every connected open set $U\subset\{x:\nabla\varphi(x)\neq0\}$, the restriction of~\eqref{eq-scalar-moment-affine-carre} to $U$ is equivalent to the existence of an eikonal coordinate $r\in C^2(U)$ and constants $c,\kappa,\lambda\in\RR$ such that
\begin{equation}\label{eq-scalar-moment-eikonal-form}
\norm{\nabla r}=1,
\qquad
\varphi=c+\kappa r+\frac{\lambda}{2}r^2.
\end{equation}
Under~\eqref{eq-scalar-moment-affine-carre}, one may take $\lambda=b/2$.
\end{prop}
\begin{proof}
The first variation is $\delta f_g(\alpha)(x)=g'(m_\varphi(\alpha))\varphi(x)$, so the continuity equation gives the displayed evolution of $m_t$. Condition~\eqref{eq-scalar-moment-affine-carre} then proves sufficiency.

For necessity, it is enough to take $g(s)=s$ and write $\psi=\norm{\nabla\varphi}^2$. Autonomy for Dirac laws first shows that $\psi=q\circ\varphi$ for some function $q$ on $\varphi(\Omega)$. Because $\Omega$ is connected, $\varphi(\Omega)$ is an interval. Comparing a Dirac law with any two-point mixture having the same $\varphi$-moment shows that
\(q\bigl((1-\zeta)s_0+\zeta s_1\bigr) = (1-\zeta)q(s_0)+\zeta q(s_1) \qquad (0\leq\zeta\leq1).\)
Hence $q$ is affine, which is precisely~\eqref{eq-scalar-moment-affine-carre}.

On a regular set $U$, the right-hand side of~\eqref{eq-scalar-moment-affine-carre} is positive. Fix $s_0\in\varphi(U)$ and define
\(r(x)=\int_{s_0}^{\varphi(x)}\frac{\d s}{\sqrt{a+bs}}.\)
Then $\norm{\nabla r}=1$. Inverting this change of variables gives~\eqref{eq-scalar-moment-eikonal-form} with $c=s_0$, $\kappa=\sqrt{a+bs_0}$, and $\lambda=b/2$. Conversely,~\eqref{eq-scalar-moment-eikonal-form} gives
$\norm{\nabla\varphi}^2=\kappa^2-2\lambda c+2\lambda\varphi$.
\end{proof}